\documentclass[11pt]{article}
\usepackage[margin=1in]{geometry}
\usepackage{amsmath,amssymb,amsthm,mathtools}
\usepackage{verbatim}
\usepackage{hyperref}
\usepackage{booktabs}

\newtheorem{theorem}{Theorem}
\newtheorem{lemma}[theorem]{Lemma}
\newtheorem{proposition}[theorem]{Proposition}
\newtheorem{corollary}[theorem]{Corollary}
\newtheorem{remark}[theorem]{Remark}
\newtheorem{conjecture}[theorem]{Conjecture}
\theoremstyle{definition}
\newtheorem{example}[theorem]{Example}

\newcommand{\Fock}{\mathcal{F}}
\newcommand{\slh}{\widehat{\mathfrak{sl}}}
\newcommand{\supp}{\operatorname{supp}}

\title{An abacus bound on the Heisenberg-crystal depth:\\
$|\kappa(\lambda)| \le k'$ for $\lambda \vdash ek'$ with empty $e$-core}
\author{Clio}
\date{2026-08-11}

\begin{document}
\maketitle

\begin{abstract}
Let $e \ge 2$ and $k' \ge 1$, and let $\kappa \colon \Pi \to \Pi$ be the
Heisenberg-crystal image map of Gerber (\cite{Gerber}, Proposition~3.12) on the
level-$1$ Fock space of $\slh_e$, computed at charge $s = 0$. We prove that for
every partition $\lambda$ of $n = ek'$ with empty $e$-core --- equivalently,
every $\lambda$ in the support of $v_{k',e} := P_e^{k'}|\emptyset\rangle$ --- the
Heisenberg-crystal depth $|\kappa(\lambda)|$ satisfies
\[
    |\kappa(\lambda)| \;\le\; k'.
\]
The proof rests on a one-line abacus observation: at level $1$, a removable
vertical $e$-strip is equivalent, on the $e$-runner abacus, to a shift of a
single bead by $e$ positions on its own runner, which decrements one part of
the $e$-quotient by exactly $1$. Combined with the empty-$e$-core hypothesis
$\sum_j |\lambda^{(j)}| = k'$, any sequence of vertical $e$-strip removals
has length at most $k'$.

We further characterize the equality case: $|\kappa(\lambda)| = k'$ if and
only if $\lambda = e \cdot \mu$ for some $\mu \vdash k'$ (i.e., each part of
$\lambda$ has multiplicity divisible by $e$). The forward direction is
proved directly from reversibility; the converse follows by combining this
with Gerber's Young-graph identification of the empty-source component
(\cite{Gerber}, Proposition~3.12(2)).

Empirically verified at six sizes $(n, e, k')\in
\{(6,3,2), (8,2,4), (9,3,3), (12,3,4), (15,3,5), (16,4,4)\}$; the equality
characterization is exact at each.
\end{abstract}

\section{Introduction and statement}

\subsection{Background}
Gerber's paper \cite{Gerber} on \emph{Heisenberg algebra, wedges and crystals}
introduces a commuting family of \emph{Heisenberg} crystal operators
$\{\widetilde b_{a,c}^{\pm 1}\}_{a \in \mathbb{Z}_{>0},\, c \in \mathbb{Z}}$
on the level-$\ell$ higher-level $q$-Fock space. At level $\ell = 1$ and charge
$s = 0$ (the setting of this note), this action lifts to a crystal on the
usual level-$1$ Fock space $\Fock$ of $\slh_e$, whose basis is indexed by
partitions.

Two structural facts about this crystal, proved in \cite[Prop.~3.12]{Gerber},
are crucial:
\begin{enumerate}
    \item[(i)] Each connected component of the H-crystal is graded and
        isomorphic (as a graded graph) to the Young graph $\mathcal{Y}$.
    \item[(ii)] The depth of a basis vector $|\lambda\rangle$ inside its
        component equals $|\kappa(\lambda)|$, where $\kappa(\lambda)$ is a
        distinguished partition (the image of $\lambda$ under the map to the
        Young graph).
\end{enumerate}
The image partition $\kappa(\lambda)$ is not readily accessible in closed
form: the constructive recipe (\cite[Thm.~5.11]{Gerber}, level-$1$
specialization) iteratively removes \emph{good vertical $e$-strips} from
$\lambda$, and $|\kappa(\lambda)|$ is the number of such removals until an
H-highest-weight vertex is reached.

\subsection{The theorem}
Write $\Pi_k$ for the set of partitions of $k$, and $\lambda^{(0)}, \ldots,
\lambda^{(e-1)}$ for the parts of the $e$-quotient of $\lambda$.

\begin{theorem}[H-crystal depth bound]
\label{thm:main}
Let $e \ge 2$ and $k' \ge 1$. For every partition $\lambda$ of $n = ek'$
with empty $e$-core,
\[
    |\kappa(\lambda)| \;\le\; k'.
\]
Moreover, equality $|\kappa(\lambda)| = k'$ holds if and only if
$\lambda = e \cdot \mu$ for some $\mu \vdash k'$, where $e \cdot \mu$ denotes
the partition in which each part of $\mu$ appears $e$ times.
\end{theorem}

\subsection{Context (post-hoc)}
This bound is the residue of a stronger attempted claim
$|\kappa(v_{k',e})| = k'$ (with $v_{k',e} = \sum_\lambda
\chi^\lambda_{(e^{k'})}\,|\lambda\rangle$), which fails category-theoretically:
Gerber's crystal is defined on basis vectors, whereas $v_{k',e}$ is a linear
combination spread across many H-components (see the note
\texttt{for-robin/2026-08-11-wake-q21-q23-downgrade.md} for detail). What
survives, and is proved here as a structural constraint on the support of
$v_{k',e}$, is the uniform depth bound.

\section{Preliminaries}

\subsection{The abacus and the $e$-quotient}
\label{sec:abacus}

Fix a partition $\lambda = (\lambda_1 \ge \lambda_2 \ge \cdots)$ and an integer
$N \ge \ell(\lambda)$ (where $\ell(\lambda)$ denotes the number of nonzero
parts). The \emph{beta-numbers of $\lambda$ with $N$ beads} are
\[
    \beta_i \;:=\; \lambda_i + (N - i), \qquad i = 1, 2, \ldots, N.
\]
These are $N$ distinct non-negative integers, strictly decreasing:
$\beta_1 > \beta_2 > \cdots > \beta_N \ge 0$.

Place \emph{beads} on the semi-infinite number line $\mathbb{Z}_{\ge 0}$ at
positions $\beta_1, \ldots, \beta_N$; call this the \emph{$e$-abacus of
$\lambda$ with $N$ beads}. Position $p \in \mathbb{Z}_{\ge 0}$ lies on
\emph{runner} $p \bmod e$, at \emph{height} $\lfloor p / e \rfloor$. Write
$N_j := |\{i : \beta_i \equiv j \pmod e\}|$ for the number of beads on runner
$j$.

\begin{remark}
Different choices of $N$ give abaci with the same $\lambda$ but differing bead
counts on each runner; the differences are bookkeeping and do not affect any
argument below.
\end{remark}

For each runner $j \in \{0, 1, \ldots, e-1\}$, let the beads on runner $j$
lie at heights $h_{j,1} > h_{j,2} > \cdots > h_{j, N_j} \ge 0$. The $j$-th
part of the \emph{$e$-quotient of $\lambda$} is the partition
$\lambda^{(j)}$ with parts
\[
    \lambda^{(j)}_i \;:=\; h_{j, i} - (N_j - i), \qquad i = 1, \ldots, N_j.
\]
(These are non-negative and weakly decreasing.) The \emph{$e$-core of
$\lambda$} is the partition obtained by pushing every bead maximally left on
its own runner; equivalently, all $\lambda^{(j)} = \emptyset$ iff $\lambda$
equals its own $e$-core iff every runner's beads occupy positions
$\{j, j+e, j+2e, \ldots, j + (N_j-1)e\}$.

We record two standard facts (proved e.g. in \cite{JamesKerber}, Chapter 2):
\begin{lemma}[Standard]
\label{lem:size-abacus}
$\displaystyle |\lambda^{(j)}| \;=\; \sum_{i=1}^{N_j} h_{j,i} \;-\; \binom{N_j}{2}.$
\end{lemma}

\begin{lemma}[Standard]
\label{lem:size-decomposition}
$\displaystyle |\lambda| \;=\; |\operatorname{core}_e(\lambda)| \;+\; e \cdot \sum_{j=0}^{e-1}
|\lambda^{(j)}|.$
\end{lemma}

\subsection{Level-$1$ vertical $e$-strips}
\label{sec:vertical-e-strips}

Rows and columns are $0$-indexed: cell $(r, c)$ sits in row $r$, column $c$.
The \emph{content} of cell $(r, c)$ is $c - r$, its residue $(c - r) \bmod e$.

Gerber's Definition~5.1 defines a \emph{vertical $e$-strip} (in the level-$\ell$
setting) as a sequence of $e$ boxes $\gamma_1, \ldots, \gamma_e$ with no
consecutive pair in the same row (``no horizontal domino''), and
$c(\gamma_{i+1}) = c(\gamma_i) - 1$ for all $i$. At level $1$ the ordering
tag on boxes is trivial; the constraint ``distinct rows plus decreasing
contents by $1$'' forces $\gamma_i = (r + i - 1, c)$ for a single common
column $c$ and starting row $r$. Thus:

\begin{proposition}[Level-$1$ vertical $e$-strip]
\label{prop:v-e-strip}
At level $1$, a vertical $e$-strip is a set of $e$ cells in a single column
at $e$ consecutive rows, i.e.,
$\{(r, c), (r+1, c), \ldots, (r+e-1, c)\}$
for some $r \ge 0$, $c \ge 0$.
\end{proposition}

A vertical $e$-strip is \emph{removable from $\lambda$} if every cell $(r+i, c)$
($0 \le i < e$) is the rightmost cell of its row in $\lambda$
(so $\lambda_{r+i+1} = c+1$ using $1$-indexed parts, i.e.\ $\lambda_{r+1} =
\cdots = \lambda_{r+e} = c+1$), and the removal (i.e., decreasing each of these
$e$ parts by $1$) yields a partition (equivalently, if row $r+e+1$ exists then
$\lambda_{r+e+1} \le c$).

\subsection{Gerber's Heisenberg crystal at level $1$}
\label{sec:H-crystal}

Following \cite[\S3.2 and Thm.~5.11]{Gerber} in the level-$1$ specialization,
the H-crystal descent operator $\widetilde b^{-1}_{1, c}$ on the basis
$\{|\lambda\rangle\}$ acts by removing the $k$-th \emph{good removable
vertical $e$-strip} of $\lambda$ (for an appropriate index $k$ selected by
the content $c$). The definition of \emph{good removable} follows the greedy
rule of \cite[Def.~5.5]{Gerber}: order the removable vertical $e$-strips by
the strict content-lex order $\succ$ (with top cell of each strip having
largest content); the good strips $X_1 \succ X_2 \succ \cdots$ are picked
greedily as $X_j$ = the $\succ$-maximum removable strip disjoint from
$X_1, \ldots, X_{j-1}$.

A basis vector $|\lambda\rangle$ is \emph{H-highest-weight} iff no vertical
$e$-strip is removable from $\lambda$ (equivalently, $\lambda$ has no $e$
consecutive equal rows). By \cite[Prop.~3.12(4)]{Gerber}, the H-depth
$|\kappa(\lambda)|$ equals the number of applications of the descent operator
(any legal sequence of good removals) required to reach an H-highest-weight
vertex.

\begin{remark}
The upper bound argument below does not require any property of the ``good''
selection: it uses only that each step is \emph{some} vertical $e$-strip removal,
and hence decrements the quotient-size sum by exactly one.
\end{remark}

\section{Key lemma: strip removal $=$ single bead shift}
\label{sec:key-lemma}

\begin{lemma}[Abacus correspondence for vertical $e$-strip removal]
\label{lem:abacus}
Let $\lambda$ be a partition with a removable vertical $e$-strip $S$ at column
$c$, rows $r, r+1, \ldots, r+e-1$. Let $\lambda' := \lambda \setminus S$. Fix
$N \ge \ell(\lambda)$; use $N$ beads for both abaci. Then the bead multiset
of $\lambda'$ is obtained from that of $\lambda$ by removing one bead at
position $c + N - r$ and adding one bead at position $c + N - r - e$. In
particular, both positions lie on runner $(c - r) \bmod e$, and the change is
a single bead shift by $e$ positions leftward on that runner.
\end{lemma}

\begin{proof}
Using $1$-indexed row indices, removability of the strip gives
$\lambda_{r+1} = \lambda_{r+2} = \cdots = \lambda_{r+e} = c + 1$, and either
$r + e = \ell(\lambda)$ or $\lambda_{r+e+1} \le c$.

The beta-numbers of $\lambda$ satisfy, for $i = 1, \ldots, e$,
\[
    \beta_{r+i} \;=\; \lambda_{r+i} + (N - r - i) \;=\; (c + 1) + (N - r - i)
    \;=\; c + N - r - i + 1.
\]
So $\beta_{r+1} = c + N - r$, $\beta_{r+2} = c + N - r - 1$, \ldots,
$\beta_{r+e} = c + N - r - e + 1$; the block
$\{\beta_{r+1}, \ldots, \beta_{r+e}\}$ is exactly the set of $e$ consecutive
integers $\{c + N - r - e + 1, \ldots, c + N - r\}$.

Removing $S$ decreases each $\lambda_{r+i}$ ($i = 1, \ldots, e$) by $1$, hence
each $\beta_{r+i}$ by $1$. The new bead positions
$\{\beta'_{r+1}, \ldots, \beta'_{r+e}\}$
are the shifted block $\{c + N - r - e, \ldots, c + N - r - 1\}$.

The symmetric difference (old block) $\triangle$ (new block) = one removal
of $c + N - r$ and one addition of $c + N - r - e$; the other $e - 1$
integers are common. All other $\beta_j$ ($j \notin \{r+1, \ldots, r+e\}$)
are unchanged.

\emph{Both positions on runner $(c - r) \bmod e$:} both differ by a multiple
of $e$, and $c + N - r \equiv c - r \pmod{e}$ (since $N$ and $N - e$ differ by
$e$; more directly $(c + N - r) - (c + N - r - e) = e$).

\emph{Position $c + N - r$ is a bead of $\lambda$:} it equals $\beta_{r+1}$
by construction.

\emph{Position $c + N - r - e$ is not a bead of $\lambda$:} we must show
$c + N - r - e \notin \{\beta_1, \ldots, \beta_N\}$. It equals
$\beta_{r+e} - 1$. In the block $\{\beta_{r+1}, \ldots, \beta_{r+e}\}$ (values
$c + N - r - e + 1, \ldots, c + N - r$), it is one below the minimum, so
absent. Above the block, $\beta_{r+e-1} = c + N - r - e + 2 > c + N - r - e$.
Below the block, either $r + e = \ell(\lambda) = N$ (no further beads with
index $> r+e$ having positive $\lambda$-part; but $\beta_j = N - j$ for
$j > \ell(\lambda)$, so $\beta_{r+e+1} = N - r - e - 1 = c + N - r - e -
(c + 1) \le c + N - r - e - 1$) or $\lambda_{r+e+1} \le c$, giving
$\beta_{r+e+1} = \lambda_{r+e+1} + (N - r - e - 1) \le c + N - r - e - 1$.
In both cases $\beta_{r+e+1} \le c + N - r - e - 1 < c + N - r - e$, and
subsequent $\beta_j$ ($j > r + e + 1$) are strictly smaller. So $c + N - r - e
\notin \{\beta_1, \ldots, \beta_N\}$, as claimed.

Combining: the bead multiset of $\lambda'$ is $(\{\beta_1, \ldots, \beta_N\}
\setminus \{c+N-r\}) \cup \{c+N-r-e\}$. This is a single bead moving from
position $c + N - r$ to $c + N - r - e$ on the runner $(c-r) \bmod e$,
leftward by $e$ positions (i.e., height decreases by $1$).
\end{proof}

\begin{corollary}[Strip removal decrements quotient-size sum by exactly $1$]
\label{cor:decrement}
Under the hypotheses of Lemma~\ref{lem:abacus}, with $j := (c - r) \bmod e$,
\[
    |\lambda'^{(j)}| \;=\; |\lambda^{(j)}| - 1, \qquad
    |\lambda'^{(j')}| \;=\; |\lambda^{(j')}| \text{ for } j' \neq j.
\]
Consequently $\sum_{i=0}^{e-1} |\lambda'^{(i)}| = \sum_{i=0}^{e-1}
|\lambda^{(i)}| - 1$, and $\operatorname{core}_e(\lambda') =
\operatorname{core}_e(\lambda)$.
\end{corollary}

\begin{proof}
By Lemma~\ref{lem:abacus}, exactly one bead on runner $j$ shifts from height
$h$ to height $h - 1$; other runners are untouched. By
Lemma~\ref{lem:size-abacus}, $|\lambda^{(j)}|$ is (sum of heights on runner
$j$) minus $\binom{N_j}{2}$; the bead-count $N_j$ is unchanged; the sum of
heights decreases by $1$. Hence $|\lambda'^{(j)}| = |\lambda^{(j)}| - 1$.

The $e$-core depends only on the bead-count profile $(N_0, \ldots, N_{e-1})$,
which is preserved by any single-runner shift. So the core is preserved.
\end{proof}

\section{Proof of Theorem~\ref{thm:main}, upper bound}
\label{sec:proof-upper}

Assume $\lambda \vdash n = ek'$ has empty $e$-core. By
Lemma~\ref{lem:size-decomposition},
\[
    ek' \;=\; |\lambda| \;=\; 0 \;+\; e \cdot \sum_{j=0}^{e-1} |\lambda^{(j)}|,
\]
so $\sum_j |\lambda^{(j)}| = k'$.

By \cite[Prop.~3.12(4)]{Gerber} (level-$1$ specialization), $|\kappa(\lambda)|$
is the length of the good-removal chain from $\lambda$ to an H-highest-weight
vertex. Denote this chain
\[
    \lambda \;=\; \lambda^{[0]} \;\to\; \lambda^{[1]} \;\to\; \cdots \;\to\;
    \lambda^{[m]},
\]
where $\lambda^{[k]} \to \lambda^{[k+1]}$ is a removal of a vertical $e$-strip
(the $\lambda^{[k+1]}$ is a specific ``good'' strip removal per Def.~5.5, but
this is not needed for the bound), $m = |\kappa(\lambda)|$, and $\lambda^{[m]}$
is H-highest-weight (no removable strips).

By Corollary~\ref{cor:decrement}, each step decrements $\sum_j |\lambda^{(j)}|$
by exactly $1$ and preserves the (empty) $e$-core. So for each $k \le m$,
\[
    \sum_{j=0}^{e-1} |\lambda^{[k], (j)}| \;=\; k' - k.
\]
Since the left side is a sum of non-negative integers, $k' - k \ge 0$, giving
$k \le k'$ for all $k$ along the chain, and in particular $m \le k'$.
This proves $|\kappa(\lambda)| \le k'$. $\square$

\section{Equality case}
\label{sec:equality}

We now identify the partitions saturating the bound. The forward direction
(``if $|\kappa| = k'$ then $\lambda = e \cdot \mu$'') is proved directly from
Lemma~\ref{lem:abacus}; the converse relies on \cite[Prop.~3.12(2)]{Gerber}.

\subsection{Forward direction}

\begin{proposition}
\label{prop:forward}
If $\lambda \vdash ek'$ has empty $e$-core and $|\kappa(\lambda)| = k'$, then
$\lambda = e \cdot \mu$ for some $\mu \vdash k'$ (i.e., every part of $\lambda$
appears with multiplicity divisible by $e$).
\end{proposition}

\begin{proof}
Suppose $|\kappa(\lambda)| = k'$. By the analysis of Section~\ref{sec:proof-upper},
the good-removal chain $\lambda = \lambda^{[0]} \to \cdots \to \lambda^{[k']}$
has length exactly $k'$, and $\sum_j |\lambda^{[k'], (j)}| = 0$. So
$\lambda^{[k']}$ has empty $e$-quotient AND empty $e$-core, hence equals
$\emptyset$ (by Lemma~\ref{lem:size-decomposition}).

Reversing the chain: $\lambda$ is obtained from $\emptyset$ by successively
\emph{adding} $k'$ vertical $e$-strips. At level $1$
(Proposition~\ref{prop:v-e-strip}) each added strip consists of $e$ cells
in a single column. Hence for every column $c \ge 0$, the number of cells of
$\lambda$ in column $c$ is $e \cdot (\text{number of strips added at column } c)$;
in particular, every column height of $\lambda$ is divisible by $e$.

Column heights of $\lambda$ are the multiplicities of parts of $\lambda'$
(the conjugate). Concretely, if the columns of $\lambda$ have heights
$e \cdot m_1 \ge e \cdot m_2 \ge \cdots \ge e \cdot m_\ell$ (with $m_i \ge 1$),
then $\lambda' = (e m_1, e m_2, \ldots, e m_\ell)$, and dually $\lambda$ has
each part appearing with multiplicity a multiple of $e$. Explicitly: if
column $c$ has height $e t_c$ ($t_c \ge 0$), then the parts of $\lambda$
equal to $j$ (for each $j \ge 1$) number $e (t_{j-1} - t_j) \ge 0$, a
non-negative multiple of $e$.

Set $\mu$ to be the partition with parts $j$ appearing with multiplicity
$t_{j-1} - t_j$; then $\lambda = e \cdot \mu$, and $|\mu| = \frac{1}{e}|\lambda|
= k'$.
\end{proof}

\subsection{Converse direction}

\begin{proposition}
\label{prop:converse}
For every $\mu \vdash k'$, $|\kappa(e \cdot \mu)| = k'$.
\end{proposition}

\begin{proof}
By \cite[Prop.~3.12(1),(2)]{Gerber}, the connected component of $|\emptyset
\rangle$ in the H-crystal is a graded, connected graph isomorphic to the
Young graph $\mathcal{Y}$; in particular, the number of vertices at
depth exactly $k'$ inside this component equals $|\Pi_{k'}|$, the number
of partitions of $k'$.

Every such depth-$k'$ vertex $|\lambda\rangle$ satisfies $|\kappa(\lambda)| =
k'$ (by \cite[Prop.~3.12(4)]{Gerber}). By Proposition~\ref{prop:forward},
each such $\lambda$ is of the form $e \cdot \mu'$ for some $\mu' \vdash k'$.
Since the map $\mu' \mapsto e \cdot \mu'$ is an injection from $\Pi_{k'}$ to
partitions, the depth-$k'$ vertices form a subset of the $|\Pi_{k'}|$-element
set $\{e \cdot \mu' : \mu' \vdash k'\}$.

But this subset has exactly $|\Pi_{k'}|$ elements, so it equals
$\{e \cdot \mu' : \mu' \vdash k'\}$. In particular, $|\kappa(e \cdot \mu)| = k'$
for every $\mu \vdash k'$.
\end{proof}

Combining Propositions~\ref{prop:forward} and~\ref{prop:converse} completes
the proof of Theorem~\ref{thm:main}.

\section{Computational verification}
\label{sec:verification}

We have implemented the level-$1$ Heisenberg-crystal descent (following
\cite[Thm.~5.11]{Gerber}: iterated good vertical $e$-strip removal) and
computed $|\kappa(\lambda)|$ for every partition $\lambda$ in $\supp(v_{k',e})
= \{\lambda \vdash ek' : \operatorname{core}_e(\lambda) = \emptyset\}$ at six
sizes.

\begin{table}[h]
\centering
\begin{tabular}{@{}lrrrl@{}}
\toprule
$(n, e, k')$ & $|\supp(v_{k',e})|$ & $\max |\kappa|$ & \#$\{\lambda:|\kappa|=k'\}$ & Distribution of $|\kappa|$ \\
\midrule
$(6, 3, 2)$   & 9   & 2 & 2 = $|\Pi_2|$ & $[5, 2, 2]$ \\
$(8, 2, 4)$   & 20  & 4 & 5 = $|\Pi_4|$ & $[5, 3, 4, 3, 5]$ \\
$(9, 3, 3)$   & 22  & 3 & 3 = $|\Pi_3|$ & $[10, 5, 4, 3]$ \\
$(12, 3, 4)$  & 51  & 4 & 5 = $|\Pi_4|$ & $[20, 10, 10, 6, 5]$ \\
$(15, 3, 5)$  & 108 & 5 & 7 = $|\Pi_5|$ & $[36, 20, 20, 15, 10, 7]$ \\
$(16, 4, 4)$  & 105 & 4 & 5 = $|\Pi_4|$ & $[51, 22, 18, 9, 5]$ \\
\bottomrule
\end{tabular}
\caption{Distribution of $|\kappa(\lambda)|$ over $\supp(v_{k',e})$. In every
case the upper bound $\max |\kappa| \le k'$ is tight, and equality is achieved
exactly by the $|\Pi_{k'}|$ partitions of the form $e \cdot \mu$.
Distribution entry $j$ = \#$\{\lambda \in \supp(v_{k',e}) : |\kappa(\lambda)|
= j\}$.}
\end{table}

The code lives in
\path{~/projects/probes/2026-08-11-q21-gerber-kappa/}
(\path{gerber_heisenberg.py}, \path{probe_extend.py}); the raw output for the
extended table is in \path{probe_extend.log}.

\section{What is (and is not) proved}
\label{sec:gaps}

\begin{itemize}
    \item \textbf{Proved rigorously here:} Theorem~\ref{thm:main} in full,
        i.e., both the bound $|\kappa(\lambda)| \le k'$ and the equality
        characterization $|\kappa(\lambda)| = k' \iff \lambda = e \cdot \mu$.
        The bound uses only the abacus correspondence
        (Lemma~\ref{lem:abacus}); the equality forward direction uses
        reversibility; the converse cites Gerber's Young-graph identification
        \cite[Prop.~3.12(2)]{Gerber}.
    \item \textbf{Cited without reproof:} the Young-graph structure of the
        empty-source H-component (Gerber Prop.~3.12(2)), and the constructive
        recipe identifying $|\kappa|$ with the good-removal chain length
        (Gerber Prop.~3.12(4) + Thm.~5.11).
    \item \textbf{Not addressed here:} the value of $\varepsilon_i(v_{k',e})$,
        Clio's Conjecture~10 (that $\varepsilon_i(v_{k',e}) = 1$ uniformly)
        and Conjecture~11 (endpoint decomposition of $\pi(w)$). This bound is
        a genuine constraint on the H-crystal geometry of the support of
        $v_{k',e}$, but does not, by itself, close either conjecture: the
        commuting-crystal argument (Def.~3.8(2)) operates on basis vectors,
        while both conjectures concern the linear combination $v_{k',e}$.
    \item \textbf{No known gap.} The proof of Theorem~\ref{thm:main} is
        complete modulo the cited external facts.
\end{itemize}

\begin{thebibliography}{9}
\bibitem{Gerber}
T.~Gerber,
\newblock \emph{Heisenberg algebra, wedges and crystals},
\newblock J. Algebraic Combin. \textbf{49}(1) (2019), 99--124.
\newblock \texttt{arXiv:1612.08760}.

\bibitem{JamesKerber}
G.~D.~James and A.~Kerber,
\newblock \emph{The Representation Theory of the Symmetric Group},
\newblock Encyclopedia of Mathematics and its Applications, Vol.~16, Addison-Wesley, 1981.
\end{thebibliography}

\end{document}
